\documentclass[10pt,a4paper]{amsart}
\usepackage{amsmath,amssymb,amsthm,graphicx}
% Retained minimal preamble for the Stacks Project chapter more-algebra.tex (More on Algebra).
% The chapter file pinned-original-source.tex carries only \input{preamble} on its first line, so no
% preamble text lives inside the chapter itself.  The Stacks Project uses one single shared preamble for
% every chapter; the snapshot stages it identically as sources/stacks-chow/preamble.tex and
% sources/stacks-etale-cohomology/preamble.tex.  The block below is copied mechanically, line for line,
% from sources/stacks-chow/preamble.tex lines 14-18 (the Xy-pic setup) and 154-217 (the native theorem
% declarations together with the mathematical macros).  Nothing was added, removed or redefined; the
% wrapper class and the amsmath/amssymb/amsthm/graphicx base packages are supplied by the assembler.
% Chapter: more-algebra.tex -- More on Algebra.
% For commutative diagrams we use Xy-pic
\usepackage[all]{xy}

% We use 2cell for 2-commutative diagrams.
\xyoption{2cell}

% Theorem environments.
%
\theoremstyle{plain}
\newtheorem{theorem}[subsection]{Theorem}
\newtheorem{proposition}[subsection]{Proposition}
\newtheorem{lemma}[subsection]{Lemma}

\theoremstyle{definition}
\newtheorem{definition}[subsection]{Definition}
\newtheorem{example}[subsection]{Example}
\newtheorem{exercise}[subsection]{Exercise}
\newtheorem{situation}[subsection]{Situation}

\theoremstyle{remark}
\newtheorem{remark}[subsection]{Remark}
\newtheorem{remarks}[subsection]{Remarks}

\numberwithin{equation}{subsection}

% Macros
%
\def\lim{\mathop{\mathrm{lim}}\nolimits}
\def\colim{\mathop{\mathrm{colim}}\nolimits}
\def\Spec{\mathop{\mathrm{Spec}}}
\def\Hom{\mathop{\mathrm{Hom}}\nolimits}
\def\Ext{\mathop{\mathrm{Ext}}\nolimits}
\def\SheafHom{\mathop{\mathcal{H}\!\mathit{om}}\nolimits}
\def\SheafExt{\mathop{\mathcal{E}\!\mathit{xt}}\nolimits}
\def\Sch{\mathit{Sch}}
\def\Mor{\mathop{\mathrm{Mor}}\nolimits}
\def\Ob{\mathop{\mathrm{Ob}}\nolimits}
\def\Sh{\mathop{\mathit{Sh}}\nolimits}
\def\NL{\mathop{N\!L}\nolimits}
\def\CH{\mathop{\mathrm{CH}}\nolimits}
\def\proetale{{pro\text{-}\acute{e}tale}}
\def\etale{{\acute{e}tale}}
\def\QCoh{\mathit{QCoh}}
\def\Ker{\mathop{\mathrm{Ker}}}
\def\Im{\mathop{\mathrm{Im}}}
\def\Coker{\mathop{\mathrm{Coker}}}
\def\Coim{\mathop{\mathrm{Coim}}}

% Boxtimes
%
\DeclareMathSymbol{\boxtimes}{\mathbin}{AMSa}{"02}

%
% Macros for moduli stacks/spaces
%
\def\QCohstack{\mathcal{QC}\!\mathit{oh}}
\def\Cohstack{\mathcal{C}\!\mathit{oh}}
\def\Spacesstack{\mathcal{S}\!\mathit{paces}}
\def\Quotfunctor{\mathrm{Quot}}
\def\Hilbfunctor{\mathrm{Hilb}}
\def\Curvesstack{\mathcal{C}\!\mathit{urves}}
\def\Polarizedstack{\mathcal{P}\!\mathit{olarized}}
\def\Complexesstack{\mathcal{C}\!\mathit{omplexes}}
% \Pic is the operator that assigns to X its picard group, usage \Pic(X)
% \Picardstack_{X/B} denotes the Picard stack of X over B
% \Picardfunctor_{X/B} denotes the Picard functor of X over B
\def\Pic{\mathop{\mathrm{Pic}}\nolimits}
\def\Picardstack{\mathcal{P}\!\mathit{ic}}
\def\Picardfunctor{\mathrm{Pic}}
\def\Deformationcategory{\mathcal{D}\!\mathit{ef}}

% --- appended 2026-10-12 (Chief): the project own definitions of the reference/slogan/history
% display environments, copied mechanically from the shared Stacks preamble lines 8-12
% (\usepackage{verbatim} plus the three \newenvironment lines, which the chapter preamble keeps
% verbatim from the project preamble). Nothing added, removed or redefined: these environments
% carry no mathematical content and the project itself typesets them as comments. Without these
% three lines a counted statement whose own text uses them fails with Environment slogan undefined.
\usepackage{verbatim}
\newenvironment{reference}{\comment}{\endcomment}
\newenvironment{slogan}{\comment}{\endcomment}
\newenvironment{history}{\comment}{\endcomment}
\usepackage{cleveref}
\crefname{originallemma}{original lemma}{original lemmas}
\Crefname{originallemma}{Original lemma}{Original lemmas}
\makeatletter
\newlabel{more-algebra-lemma-locally-nilpotent-henselian}{{0ALI}{source}{}{sourcepointer-1}{}}
\newlabel{more-algebra-lemma-locally-nilpotent-henselian@cref}{{[originallemma][0][]0ALI}{[1][0][]source}{}{}{}}
\makeatother
\crefname{originallemma}{original lemma}{original lemmas}
\Crefname{originallemma}{Original lemma}{Original lemmas}
\makeatletter
\newlabel{algebra-lemma-lift-etale}{{04D1}{source}{}{sourcepointer-2}{}}
\newlabel{algebra-lemma-lift-etale@cref}{{[originallemma][0][]04D1}{[1][0][]source}{}{}{}}
\makeatother
\crefname{originallemma}{original lemma}{original lemmas}
\Crefname{originallemma}{Original lemma}{Original lemmas}
\makeatletter
\newlabel{algebra-lemma-smooth-strong-lift}{{07K4}{source}{}{sourcepointer-3}{}}
\newlabel{algebra-lemma-smooth-strong-lift@cref}{{[originallemma][0][]07K4}{[1][0][]source}{}{}{}}
\makeatother
\crefname{originallemma}{original lemma}{original lemmas}
\Crefname{originallemma}{Original lemma}{Original lemmas}
\makeatletter
\newlabel{algebra-lemma-diagonal-unramified}{{02FL}{source}{}{sourcepointer-4}{}}
\newlabel{algebra-lemma-diagonal-unramified@cref}{{[originallemma][0][]02FL}{[1][0][]source}{}{}{}}
\makeatother
\crefname{originallemma}{original lemma}{original lemmas}
\Crefname{originallemma}{Original lemma}{Original lemmas}
\makeatletter
\newlabel{algebra-lemma-unramified}{{00UV}{source}{}{sourcepointer-5}{}}
\newlabel{algebra-lemma-unramified@cref}{{[originallemma][0][]00UV}{[1][0][]source}{}{}{}}
\makeatother
\crefname{originallemma}{original lemma}{original lemmas}
\Crefname{originallemma}{Original lemma}{Original lemmas}
\makeatletter
\newlabel{algebra-lemma-locally-nilpotent}{{0544}{source}{}{sourcepointer-6}{}}
\newlabel{algebra-lemma-locally-nilpotent@cref}{{[originallemma][0][]0544}{[1][0][]source}{}{}{}}
\makeatother
\crefname{originallemma}{original lemma}{original lemmas}
\Crefname{originallemma}{Original lemma}{Original lemmas}
\makeatletter
\newlabel{algebra-lemma-lift-idempotents}{{00J9}{source}{}{sourcepointer-7}{}}
\newlabel{algebra-lemma-lift-idempotents@cref}{{[originallemma][0][]00J9}{[1][0][]source}{}{}{}}
\makeatother
\crefname{originallemma}{original lemma}{original lemmas}
\Crefname{originallemma}{Original lemma}{Original lemmas}
\makeatletter
\newlabel{lemma-lift-factorization-monic}{{0ALH}{source}{}{sourcepointer-8}{}}
\newlabel{lemma-lift-factorization-monic@cref}{{[originallemma][0][]0ALH}{[1][0][]source}{}{}{}}
\makeatother
\title{Pairs with locally nilpotent ideal are henselian and etale algebras reduce bijectively}
\author{The Stacks Project Authors}
\date{stacks-project:more-algebra.tex:local-snapshot-20261012; GNU FDL 1.2 or later}
\begin{document}
\maketitle
\setcounter{subsection}{1}
\noindent\textit{Editorial note.} All delivered mathematical text is the author's own, copied from the pinned Stacks Project chapter file more-algebra.tex; no statement or proof step was rewritten, reordered or removed.\par\medskip
\noindent\textit{Editorial note.} The chapter file pinned-original-source.tex carries only \textbackslash{}input\{preamble\} on its first line, so no preamble text lives inside the chapter. The retained file sources/stacks-more-algebra/retained-preamble.tex is a mechanical, line-for-line copy of the project's single shared preamble as staged in the snapshot (Xy-pic setup plus theorem declarations and mathematical macros). Nothing was added, removed or redefined, and the excerpt serves only as preamble, never as a span coordinate.\par\medskip
\noindent\textit{Editorial note.} Both counted spans start at the \textbackslash{}begin line of their environment, so the statement and the author's proof are delivered whole; surrounding prose of the section is deliberately outside the counted spans and is carried as standard prerequisites instead.\par\medskip
\noindent\textit{Editorial note.} The citation fields (source\_id, source\_version, authors, source\_url, license.name, source\_title) are identical to those of every other card this lane has produced from this chapter, so the batch selector treats them as one source.\par\medskip
\noindent\textit{Editorial note.} Benign-screen note: the proof contains the phrase 'it is clear that' / 'is obvious' once, used as a legitimate step inside a complete argument; the author proves the statement, so this is not a hint-only proof.\par\medskip
\noindent\textit{Allowed standard prerequisites (editorial).}\begin{itemize}
\item Lifting of etale algebras, and the strong lifting property for smooth algebras. Algebra, Lemma \textbackslash{}ref\{algebra-lemma-lift-etale\} and Algebra, Lemma \textbackslash{}ref\{algebra-lemma-smooth-strong-lift\}, cited at proof lines 2424 and 2428-2429. Give essential surjectivity and fullness.
\item An etale algebra is unramified, and the kernel of the multiplication map is generated by an idempotent. Algebra, Lemma \textbackslash{}ref\{algebra-lemma-diagonal-unramified\} and Algebra, Lemma \textbackslash{}ref\{algebra-lemma-unramified\}, cited at proof lines 2433-2435. Provide the idempotent used for faithfulness.
\item Idempotents lift modulo a locally nilpotent ideal, and the ideal is locally nilpotent after base change. Algebra, Lemma \textbackslash{}ref\{algebra-lemma-lift-idempotents\} and Algebra, Lemma \textbackslash{}ref\{algebra-lemma-locally-nilpotent\}, cited at proof lines 2437-2438. Give faithfulness and support the radical argument.
\item Lifting of a factorisation of a monic polynomial into monic factors after an etale base change inducing an isomorphism modulo the ideal. Lemma \textbackslash{}ref\{lemma-lift-factorization-monic\}, cited at proof lines 2445-2447. Provides the factorisation used in the henselian verification.
\end{itemize}\medskip
\ifdefined\hypertarget\hypertarget{sourcepointer-1}{}\fi
\noindent\textit{Source reference.} Original lemma 0ALI: focused-seeds/AGP/sources/stacks-snapshot-20261011/tags.txt, source lines 13371--13371. The lemma's own canonical tag, quoted so that readers can resolve the label inside the delivered text; it points at the counted statement itself and carries no proof content.\par\medskip
\ifdefined\hypertarget\hypertarget{sourcepointer-2}{}\fi
\noindent\textit{Source reference.} Original lemma 04D1: focused-seeds/AGP/sources/stacks-snapshot-20261011/tags.txt, source lines 5711--5711. Navigation pointer to a statement the counted proof refers to but does not contain; nothing load-bearing is delegated to it.\par\medskip
\ifdefined\hypertarget\hypertarget{sourcepointer-3}{}\fi
\noindent\textit{Source reference.} Original lemma 07K4: focused-seeds/AGP/sources/stacks-snapshot-20261011/tags.txt, source lines 9644--9644. Navigation pointer to a statement the counted proof refers to but does not contain; nothing load-bearing is delegated to it.\par\medskip
\ifdefined\hypertarget\hypertarget{sourcepointer-4}{}\fi
\noindent\textit{Source reference.} Original lemma 02FL: focused-seeds/AGP/sources/stacks-snapshot-20261011/tags.txt, source lines 3196--3196. Navigation pointer to a statement the counted proof refers to but does not contain; nothing load-bearing is delegated to it.\par\medskip
\ifdefined\hypertarget\hypertarget{sourcepointer-5}{}\fi
\noindent\textit{Source reference.} Original lemma 00UV: focused-seeds/AGP/sources/stacks-snapshot-20261011/tags.txt, source lines 1125--1125. Navigation pointer to a statement the counted proof refers to but does not contain; nothing load-bearing is delegated to it.\par\medskip
\ifdefined\hypertarget\hypertarget{sourcepointer-6}{}\fi
\noindent\textit{Source reference.} Original lemma 0544: focused-seeds/AGP/sources/stacks-snapshot-20261011/tags.txt, source lines 6627--6627. Navigation pointer to a statement the counted proof refers to but does not contain; nothing load-bearing is delegated to it.\par\medskip
\ifdefined\hypertarget\hypertarget{sourcepointer-7}{}\fi
\noindent\textit{Source reference.} Original lemma 00J9: focused-seeds/AGP/sources/stacks-snapshot-20261011/tags.txt, source lines 707--707. Navigation pointer to a statement the counted proof refers to but does not contain; nothing load-bearing is delegated to it.\par\medskip
\ifdefined\hypertarget\hypertarget{sourcepointer-8}{}\fi
\noindent\textit{Source reference.} Original lemma 0ALH: focused-seeds/AGP/sources/stacks-snapshot-20261011/tags.txt, source lines 13370--13370. Navigation pointer to a statement the counted proof refers to but does not contain; nothing load-bearing is delegated to it.\par\medskip
\begin{lemma}
\label{lemma-locally-nilpotent-henselian}
Let $(A, I)$ be a pair with $I$ locally nilpotent. Then the functor
$B \mapsto B/IB$ induces an equivalence between the category of
\'etale algebras over $A$ and the category of \'etale algebras over $A/I$.
Moreover, the pair is henselian.
\end{lemma}

\begin{proof}
Essential surjectivity holds by Algebra, Lemma \ref{algebra-lemma-lift-etale}.
If $B$, $B'$ are \'etale over $A$ and $B/IB \to B'/IB'$ is a morphism
of $A/I$-algebras, then we can lift this by
Algebra, Lemma \ref{algebra-lemma-smooth-strong-lift}.
Finally, suppose that $f, g : B \to B'$ are two $A$-algebra
maps with $f \bmod I = g \bmod I$. Choose an idempotent $e \in B \otimes_A B$
generating the kernel of the multiplication map $B \otimes_A B \to B$,
see Algebra, Lemmas \ref{algebra-lemma-diagonal-unramified}
and \ref{algebra-lemma-unramified} (to see that \'etale is unramified).
Then $(f \otimes g)(e) \in IB'$. Since $IB'$ is locally nilpotent
(Algebra, Lemma \ref{algebra-lemma-locally-nilpotent}) this implies
$(f \otimes g)(e) = 0$ by Algebra, Lemma \ref{algebra-lemma-lift-idempotents}.
Thus $f = g$.

\medskip\noindent
It is clear that $I$ is contained in the Jacobson radical of $A$.
Let $f \in A[T]$ be a monic polynomial and let
$\overline{f} = g_0h_0$ be a factorization
of $\overline{f} = f \bmod I$ with $g_0, h_0 \in A/I[T]$ monic
generating the unit ideal in $A/I[T]$. By
Lemma \ref{lemma-lift-factorization-monic}
there exists an \'etale ring map $A \to A'$ which
induces an isomorphism $A/I \to A'/IA'$ such that
the factorization lifts to a factorization into monic polynomials
over $A'$. By the above we have $A = A'$ and the factorization
is over $A$.
\end{proof}

\end{document}
